\chapter{Conclusion and Future Work}
\label{ch:conclusion}

\section{Conclusion}
\label{sec:conclusion}

We have described the Sprout conversational AI subsystem: the architecture, algorithms, and implementation that allow the platform to communicate with a learner across four modalities, record interactions durably, and compose them into adaptive teaching. Replacing a multi-service Python/LiveKit pipeline with a thin, JWT-validated WebSocket proxy fronting Gemini Live has produced lower dialogue latency, operational complexity, and infrastructure cost than the legacy baseline. Six algorithms make this concrete; an explicit Firestore schema turns the interaction stream into structured information for dashboards, billing, and audit.

The forward trajectory enriches the conversational layer with tools rather than more modalities. Evaluation will strengthen through the online protocol in Section~\ref{sec:online-protocol}; personalization will strengthen through the multi-intent representation in Section~\ref{sec:multi-intent}.

\section{Function Calling}
\label{sec:function-calling}

Function calling allows the model to emit structured calls into a registry of caller-defined functions~\cite{sprout2026function}. A Chrome-extension prototype demonstrated:
\begin{itemize}
    \item \texttt{show\_hint(coords, color, duration)} --- visual indicator on a screen grid.
    \item \texttt{display\_image(query)} --- educational image retrieval or generation.
    \item \texttt{navigate\_to\_website(url, returnStrategy)} --- curated external resource in a side pane.
\end{itemize}

The proxy already routes function-call frames distinctly in Algorithm~\ref{alg:a2}. Implementation requires specialist services, a function registry, and typed result frames.

\section{Multi-Agent Coordination}
\label{sec:multi-agent}

A larger shift decomposes the single Gemini Live session into a coordinator--specialist system~\cite{midscene2024,sprout2026multiagent}:
\begin{itemize}
    \item \textbf{Visual Learning Agent} --- \texttt{display\_image} and visual tooling.
    \item \textbf{Navigation Agent} --- \texttt{navigate\_to\_website} and return-strategy logic.
    \item \textbf{Assessment Agent} --- lightweight skill probes without exiting the conversation.
    \item \textbf{Accessibility Agent} --- pacing and modality adaptation for neurodiverse learners.
\end{itemize}

A central Learning Coordinator maintains \texttt{LearningContext} (profile, session state, knowledge-map node, recent turns) and selects specialists by capability, recent performance, and cost. Adding capabilities is additive: new specialists register new functions without changing the proxy transport.

\section{Calibrated Lesson Recommendation}
\label{sec:calibrated-rec}

Calibrated lesson recommendation draws on diversification literature~\cite{steck2018calibrated,carbonell1998mmr,anh2026mmr}. After several lesson completions, the system recommends from eligible knowledge-map successors using a calibrated reranking step that mirrors the empirical distribution of the learner's interests weighted by recency.

\textbf{Relationship to Section~\ref{sec:multi-intent}.} Calibrated reranking operates at \emph{lesson selection} time within a course; \textsc{BuildSessionContext} operates at \emph{session priming} time across courses. Together they address multi-interest learners: priming supplies distinct intent vectors per active course (Section~\ref{sec:multi-intent}); reranking ensures the next lesson within a course reflects calibrated diversity rather than greedy relevance maximisation. Neither replaces the other; both should be evaluated in the Phase~3 A/B protocol (Section~\ref{sec:online-protocol}).

\section{On-Device Privacy Processing}
\label{sec:on-device}

Webcam frames currently flow to Gemini Live. For posture or eye-contact features, on-device models could extract low-dimensional summaries (posture score, gaze flag) so only summaries leave the device, reducing exposure and bandwidth.

\section{Closing Remarks}
\label{sec:closing}

The system documented here is a small, opinionated foundation. It bets that a thin proxy in front of a managed multimodal model is the right abstraction for an educational tutor---and constructs a typed SDK, explicit Firestore schema, six algorithms, observability stack, a multi-intent priming extension, and a rigorous evaluation roadmap around that bet. Pilot measurements and reduced operational complexity suggest the bet has paid off; function calling, multi-agent coordination, and online experimentation are the additive next steps.
